\subsection{性健康 App 与可穿戴设备：实用指南与陷阱}

智能手机与可穿戴设备的普及，让"性健康"第一次可以被普通人日常量化与自我管理。周期追踪、盆底训练、勃起监测、亲密关系打卡……这些工具既有真实价值，也暗藏误区。本节从"能做什么、证据如何、有什么坑"三个角度，给出一份实用指南。

\subsubsection{常见类型与用途}

\begin{itemize}
	\item \textbf{周期与生育追踪}：记录月经、排卵、基础体温与宫颈黏液，辅助备孕或自然避孕。其核心逻辑是"症状体温法（Symptothermal Method）"。但请注意：仅凭 App 算法推算的"安全期"避孕失败率远高于规范使用的激素或屏障避孕，\textbf{不可作为唯一的避孕手段}。
	\item \textbf{盆底训练（凯格尔）App}：通过定时提示与生物反馈，帮助建立正确的盆底收缩。对压力性尿失禁与产后盆底康复有中等质量证据支持；但\textbf{"练错肌肉"极为常见}，最好先经专业人员评估。
	\item \textbf{性功能自评量表}：如男性国际勃起功能指数（IIEF-5）、女性性功能指数（FSFI）。量表有临床参考价值，但\textbf{自评不等于诊断}，分数偏低应就医，而非自行用药。
	\item \textbf{亲密关系与沟通类}：伴侣问答、欲望差异协商、性爱日历等，本质是把"沟通"结构化，可作为关系对话的破冰工具。
	\item \textbf{性教育与咨询}：提供科普内容与在线问诊入口，便利但有质量参差，应认准有资质的机构与实名医生。
\end{itemize}

\subsubsection{可穿戴设备能测什么}

\begin{itemize}
	\item \textbf{睡眠、心率变异性（HRV）与体温}：手环与智能戒指可连续记录这些指标。它们与性欲、勃起和整体状态的关联是\textbf{间接的}——长期睡眠不足、HRV 持续偏低、体温节律紊乱，往往伴随性欲下降；但设备本身\textbf{并不能直接"测性欲"}。
	\item \textbf{夜间阴茎胀大试验（NPT）}：这是临床上区分\textbf{心理性}与\textbf{器质性}勃起功能障碍的重要工具。传统方法包括\textbf{邮票试验}（睡前环绕阴茎贴邮票，晨起查看是否断裂）与 \textbf{RigiScan} 等设备监测。若夜间与晨间勃起良好、而清醒时勃起困难，多提示以心理因素为主。此类监测应在医生指导下进行。
	\item \textbf{商业"勃起监测"与性健康手环}：不少产品宣称可"评估性功能""提升性欲"，但\textbf{多数缺乏临床验证}，其算法与"正常值"往往不透明。设备数据可作参考，不能替代医学检查。
\end{itemize}

\subsubsection{必须警惕的三个陷阱}

\begin{itemize}
	\item \textbf{证据缺口}：绝大多数性健康 App \textbf{未经临床验证}，其"评分""预测"多为经验或商业设定，不能等同于医疗诊断。
	\item \textbf{数据隐私}：性健康数据是\textbf{最敏感的个人信息之一}，一旦泄露后果严重。使用前应查看隐私政策，优先选择\textbf{本地存储、可导出、可删除}且数据安全承诺明确的产品；谨慎授权通讯录、定位与云端同步。
	\item \textbf{过度量化焦虑}：把亲密关系变成每天打分的"KPI"，可能反过来\textbf{加重表现焦虑}、损害自然的性反应。工具应服务于人，而不是让人被数字绑架。
\end{itemize}

\begin{tcolorbox}[colback=blue!5!white,colframe=blue!50!black,title=贴心小叮咛]
App 与可穿戴设备是\textbf{自我觉察的辅助工具}，不是医生。它们擅长"记录与提醒"，不擅长"诊断与治疗"。若出现持续勃起困难、性欲显著下降、性交疼痛或伴侣关系困扰，\textbf{请直接就医}，别把希望寄托在评分与"智能建议"上。花钱买设备之前，先问一句：它解决的是我的问题，还是制造了新问题？
\end{tcolorbox}
